\section{Limitations}
\label{sec:limits}

\subsection{Run-to-run variance and multi-seed dispersion}

All reported dispersions are the sample standard deviation over three independent
seeds (one run per seed, \texttt{ddof}=1): \PublishedSdAuc{} on AUC and \PublishedSdRecall{} on
recall across the tables. Two caveats follow. First, this conflates between-seed variance with
run-to-run variance; the run-to-run spread on the \emph{same} stream was not re-measured after
the generator de-leakage, so no run-to-run figure is reported. Second, with three seeds, the
paired Wilcoxon signed-rank test has a minimum attainable $p$ of $0.25$, so the configuration
node's lateral-improvement deltas (\BlatAucD{} AUC, \BlatRecD{} routed recall) cannot be called
statistically significant as such; the sign pattern is reported instead (positive on
all three seeds), and the effect is treated as consistent but unproven. Extending evaluations to $\geq$5
seeds with $\geq$3 same-stream replicates remains future work.

\subsection{Synthetic benchmark fidelity}

Synthetic access streams cannot capture the full behavioral diversity of production
enterprise networks. We constrain artificial separability through the automated
leakage audit (Section~\ref{sec:setup}) and ground the 5-node schema in real enterprise
telemetry via PicoDomain (Section~\ref{sec:external}).

\subsection{Operating recall trade-offs}

Cost-sensitive threshold routing maps ranking discriminability to concrete operating points,
yet subtle lateral movement partially overlaps with legitimate exploratory access.
Operating thresholds must therefore be calibrated against SOC triage capacity and
acceptable alert volume rather than unweighted ranking metrics.

\subsection{State complexity and scalability}
\label{sec:limits-state}

Constant-time tensor lookups and bounded memory consumption apply directly to the GPU/CPU
tensor representations (entity memory vectors, neighbor ring buffers, and static feature tables),
which are allocated at fixed capacity. However, peripheral structures such as interaction frequency
counters and precursor alert trackers scale with the volume of distinct entity pairs. Eviction
mechanisms require buffer scanning for invalidation. While measured serving latency remains within
operational limits ($P_{50} \approx \LatencyPFifty$\,ms, $P_{99} \approx \LatencyPNinetyNine$\,ms),
deployments serving millions of concurrent entities require distributed cache eviction strategies.

\subsection{Commit gate boundary conditions}

As detailed in Section~\ref{sec:problem}, external commit gating admits authorized
adversary actions into entity memory while risking update starvation for persistently
flagged benign entities. Resolving these boundary conditions requires operational
grace periods and periodic administrative overrides.

\subsection{Precursor heuristic scope}

The precursor prior operates on the assumption that lateral movement often follows detectable
reconnaissance activity on the same origin host. An adversary conducting silent reconnaissance
or moving directly from a pre-established foothold will not trigger the additive precursor prior,
leaving detection to the graph structural and feature heads alone.

\subsection{Cold start dynamics}

Newly introduced entities lack sufficient historical interactions for the model to establish
habitual baselines. While evaluation streams allow entities to warm up before measuring lateral
movement, enterprise environments exhibit continuous entity churn. During initial interaction
phases, detection sensitivity is reduced until minimum interaction counts are established.

\subsection{Fingerprint evasion}

The client-configuration node detects unfamiliar client tool stacks associated with known
principals. If an adversary faithfully replicates the client TLS parameters,
ciphersuite ordering, and client environment, the request appears structurally habitual on the
\textit{config}\,$\rightarrow$\,\textit{user} binding. The configuration node increases adversary
work factor, but does not prevent exact client cloning.

\subsection{Service deployment security}

The serving prototype does not enforce cryptographic mutual TLS or API token authentication
at the transport layer. In production deployments, inference and commit endpoints must be
restricted to internal network fabrics with mutual authentication to prevent unauthorized
state manipulation.
